\section{Network moment laws: what PLL tuning can and cannot change}
\label{sec:networkmoments}

The preceding interaction construction protects finite-frequency modal decay.
We now ask a different question: which features of the low-frequency response
are reachable by PLL tuning at all? This separates controller authority from
an attractive but potentially misleading interpretation as adjustable inertia.
The derivation uses the actual lossy network and device dynamics. It does not
replace their Jacobian by a symmetric physical Laplacian. The simple-pole
resolvent and moment identities are classical \cite{beyn,malladamoments};
the model-specific content is their identification with the delayed PLL ports.

\subsection{All-PLL elimination and assumptions}
Hold replacement, dispatch, physical plant and operating point fixed. Remove
each PLL triple $(\theta_i,\omega_i,\xi_i)$ from the hidden state, retaining
the other device and network variables in $z$. The linearized port equations are
\begin{align}
 (sI-A_h)z&=B_\theta\theta+B_d d,\\
 e&=C_hz+D_\theta\theta+D_dd,\\
 \phi&=O_hz+O_\theta\theta+O_dd.
\end{align}
Here $d$ is a disturbance parameter, $\phi$ contains bus voltage phases, and
the direct terms include the event-dependent algebraic voltage response.
Assume $-A_h$ is invertible. Define the analytic, gain-independent maps
\begin{align}
 G&=D_\theta+C_h(sI-A_h)^{-1}B_\theta,&
 g&=D_d+C_h(sI-A_h)^{-1}B_d,\\
 T&=O_\theta+O_h(sI-A_h)^{-1}B_\theta,&
 Z&=O_d+O_h(sI-A_h)^{-1}B_d.
\end{align}
The actual filtered type-II PLL gives
\[
 s^2(1+t_{f,i}s)\theta_i=(sK_{p,i}+K_{i,i})e^{-s\tau_i}e_i.
\]
For nonzero integral gains this yields the exact retained equation
\begin{equation}
 K(s)\theta=g(s)d,\qquad
 K(s)=\diag\!\left(
 \frac{s^2(1+t_{f,i}s)e^{s\tau_i}}{sK_{p,i}+K_{i,i}}
 \right)-G(s).
 \label{eq:momentK}
\end{equation}
It is valid near zero, away from singular eliminated blocks. It is not a
global pole-equivalence claim across those singularities.

Write $G(s)=\sum_j s^jG_j$, and similarly for the other analytic maps.
The coefficients, not derivatives without factorials, are
\begin{align}
 K_0&=-G_0,& K_1&=-G_1,\\
 K_2&=-G_2+\diag(K_i^{-1}),&
 K_3&=-G_3+\diag\!\left((t_f+\tau)/K_i-K_p/K_i^2\right).
 \label{eq:momentKcoeff}
\end{align}
Rotational equivariance supplies $K_0r=0$ with $r=\one$, and $T_0r=\one$
for physical bus phases. Additionally require
\[
 \ker K_0=\operatorname{span}\{r\},\quad
 \ell^*K_0=0,\quad d_0=\ell^*K_1r\ne0.
\]
Then the classical simple-root residue is
\begin{equation}
 K(s)^{-1}=\frac{P}{s}+O(1),\qquad
 P=\frac{r\ell^*}{d_0}.
 \label{eq:momentresidue}
\end{equation}
The gauge is retained for this Laurent calculation and is deflated separately
for physical stability tests. Merely deleting a coordinate before the
calculation would discard the residue responsible for the frequency response.

For a common measurement window $W>0$, the Hz output transfer is
\begin{equation}
 H(s)=\frac{1-e^{-sW}}{2\pi W}
 \left[T(s)K(s)^{-1}g(s)+Z(s)\right]
      =H_0+sH_1+s^2H_2+O(s^3).
 \label{eq:momentfreq}
\end{equation}
Its removable singularity gives $H_0=T_0Pg_0/(2\pi)$. The instantaneous
phase-derivative output follows by taking $W\to0$.

\status{EXACT\_IDENTITY}
\begin{theorem}[Finite PLL changes and the synchronization residue]
\label{thm:momentlaw}
Under the assumptions above, at fixed plant and replacement:
\begin{enumerate}
\item $H_0$ is independent of all PLL proportional gains, nonzero integral
gains, and fixed measurement delays.
\item For any finite integral-gain change, allowing proportional gains and
delays to change as well, let
\[
 D_2=\diag(1/K_i^{b}-1/K_i^{a}),\qquad
 \kappa_2=\ell^*D_2r/d_0.
\]
Then $H_b-H_a=-s\kappa_2H_0+O(s^2)$.
\item If integral gains are held fixed, let
\[
 D_3=\diag(\Delta\tau/K_i-\Delta K_p/K_i^2),\qquad
 \kappa_3=\ell^*D_3r/d_0.
\]
Then
\begin{equation}
 \boxed{H_b(s)-H_a(s)=-s^2\kappa_3H_0+O(s^3).}
 \label{eq:momentlaw}
\end{equation}
These are exact identities between Taylor coefficients of the full linearized
model. They are not finite-frequency approximations with a certified error
range, nor identities for finite-amplitude nonlinear trajectories.
\end{enumerate}
\end{theorem}
\begin{proof}
The first two $K$ coefficients and hence the residue are independent of PLL
gains and delays. For two designs the exact resolvent identity gives
\[
 K_b^{-1}-K_a^{-1}=-K_b^{-1}(K_b-K_a)K_a^{-1}.
\]
If $K_b-K_a=s^mD_m+O(s^{m+1})$, $m\ge2$, substitution of
\eqref{eq:momentresidue} gives
\[
 K_b^{-1}-K_a^{-1}=-s^{m-2}PD_mP+O(s^{m-1}).
\]
Since $PD_mP=(\ell^*D_mr/d_0)P$ and the window factor is
$s/(2\pi)+O(s^2)$, the frequency difference is
$-s^{m-1}(\ell^*D_mr/d_0)H_0+O(s^m)$.
Equation~\eqref{eq:momentKcoeff} identifies $m=2$ or $m=3$ and its finite
coefficient difference. No small parameter change is required.
\end{proof}

\subsection{An exact restriction on frequency-response shaping}
Suppose both nongauge systems are exponentially stable and the relevant
step-response moments converge. For a common step $d(t)=d_0^{\rm in}$,
write $f_\infty=H_0d_0^{\rm in}$. The Laplace-transform moment formula gives
\begin{align}
 \int_0^\infty(f_b-f_a)\,dt&=-\kappa_2f_\infty
       &&\text{if integral gains change},\label{eq:arealaw}\\
 \int_0^\infty(f_b-f_a)\,dt&=0
       &&\text{if integral gains are fixed},\\
 \int_0^\infty t(f_b-f_a)\,dt&=\kappa_3f_\infty
       &&\text{if integral gains are fixed}.\label{eq:firstmomentlaw}
\end{align}
These are signed integrals. Cancellation permits large excursions and does
not bound nadir, absolute-error integrals, RoCoF, or settling time.
If a component of $f_\infty$ vanishes, normalizing by it is invalid.

\begin{corollary}[One direction of first-area authority]
For nonzero final response define
$A_j=\int_0^\infty(1-f_j(t)/f_{\infty,j})\,dt$.
Any admissible finite PLL gain/delay change at fixed plant satisfies
\[
 A_j^b-A_j^a=\kappa_2,\qquad
 (A_j-A_k)^b=(A_j-A_k)^a.
\]
Thus PLL tuning cannot independently change the buswise normalized signed
areas. Proportional gains and delays cannot change any of these areas at all.
\end{corollary}
This is a limitation of the declared architecture, not a theorem that GFLs
cannot support frequency. Altering active-power control, synthetic-inertia
feedthrough, droop, dispatch or replacement changes the premise. In particular,
an SG-to-GFL change modifies $G,g,T,Z$ and invalidates reuse of the same weights.

\subsection{Signed network weights and spatial latency}
Define the real weights
\begin{equation}
 w_i=\bar\ell_i r_i/d_0,\qquad
 \kappa_3=\sum_iw_i\left(\Delta\tau_i/K_{i,i}
                          -\Delta K_{p,i}/K_{i,i}^2\right).
 \label{eq:networkweights}
\end{equation}
They incorporate the actual network, device dynamics, operating point and
replacement. They are not nonnegative inertia shares, probabilities, or a
Laplacian centrality. Negative weights allow the leading signed moment to
respond with either sign to an equal positive delay increment at different
sites. That sign alone says nothing about stability or frequency peaks.

For a prescribed delay-increment multiset $t_1\le\cdots\le t_n$, put
$a_i=w_i/K_{i,i}$ and order $a_{(1)}\le\cdots\le a_{(n)}$.
The rearrangement inequality gives the exact extreme \emph{moment} coefficients
over all assignments:
\[
 \min_\pi\sum_i a_it_{\pi(i)}=\sum_i a_{(i)}t_{n+1-i},\qquad
 \max_\pi\sum_i a_it_{\pi(i)}=\sum_i a_{(i)}t_i.
\]
This is a closed-form placement law for a response coefficient, not an
optimization of physical delays or of the replacement frontier. A graph
Fourier expansion of $w$ merely changes coordinates unless its modal content
is independently shown to explain a design consequence.

\subsection{Collective compensation and competing modal authority}
Keeping $K_i$ fixed, preserve $H_2$ by imposing the single exact equality
\begin{equation}
 c^Tk=b_\tau,\qquad
 c_i=w_i/K_{i,i}^2,\quad
 b_\tau=\sum_iw_i\Delta\tau_i/K_{i,i},\quad k=\Delta K_p.
 \label{eq:collectivemoment}
\end{equation}
The sitewise rule $k_i=K_{i,i}\Delta\tau_i$ is sufficient but unnecessary.
For box bounds $l\le k\le u$, the equality is feasible exactly when
\[
 \sum_i\min(c_il_i,c_iu_i)\le b_\tau\le
 \sum_i\max(c_il_i,c_iu_i).
\]
This proves feasibility only of moment compensation. Without box bounds,
minimizing $k^TRk/2$ gives the standard projection
\[
 k_0=\frac{R^{-1}c}{c^TR^{-1}c}b_\tau.
\]
Suppose a simple critical root has local gain gradient $a=\nabla_{K_p}\Re\lambda$
and requires $a^Tk\le h$. The remaining damping direction is
\[
 v=R^{-1}a-R^{-1}c\frac{c^TR^{-1}a}{c^TR^{-1}c}.
\]
If $a^Tk_0>h$ and $a^Tv>0$, the exact solution of this \emph{linearized}
one-mode subproblem is
\begin{equation}
 k^*=k_0-v\frac{a^Tk_0-h}{a^Tv},\qquad
 \frac12{k^*}^TRk^*
 =\frac{b_\tau^2}{2c^TR^{-1}c}
  +\frac{(a^Tk_0-h)^2}{2a^Tv}.
 \label{eq:projectedauthority}
\end{equation}
If $a^Tv=0$ and the inequality fails at $k_0$, no moment-neutral gain direction
can satisfy that local modal requirement. For multiple modes or active boxes,
solve the corresponding convex QP and re-evaluate the exact DDE roots.
This is standard constrained optimization, not a new global optimality theorem.

The root sensitivities use the full exponential characteristic
\[
 \lambda_p=-\frac{y^*\Delta_p(\lambda)x}
                  {y^*\Delta_s(\lambda)x},\quad
 \Delta(s)=sI-A_0-B\diag(e^{-s\tau})C.
\]
Root tracking, line-search/globalization and complete-spectrum checks are
separate requirements. A one-mode correction can transfer the active limit to
another mode. Moments cannot detect that failure.

\subsection{Graph-polynomial actions: a conditional design direction}
Consider an explicitly different architecture with invertible matrix integral
gain $K_I$ and matrix proportional gain $K_P$, acting on delayed detector
errors. Its PLL equation is
\[
 s^2(I+sT_f)\theta=(sK_P+K_I)e^{-sT_\tau}e.
\]
The retained controller coefficients are now
\[
 K_2^{\rm ctrl}=K_I^{-1},\qquad
 K_3^{\rm ctrl}=T_\tau K_I^{-1}+K_I^{-1}T_f
                         -K_I^{-1}K_PK_I^{-1}.
\]
The order and residue theorem still applies; matrix ordering matters.
At fixed $K_I$ and delays, choose
\begin{equation}
 \Delta K_P=K_I S K_I,\qquad Sr=0.
 \label{eq:graphneutral}
\end{equation}
Then $\Delta K_3^{\rm ctrl}r=-Sr=0$, so $H_0,H_1,H_2$ are unchanged exactly.
For a valid graph Laplacian $L$ with $Lr=0$, any polynomial
$S=\sum_{j\ge1}\beta_jL^j$ has this property. Such an action supplies
finite-frequency tuning directions that are automatically neutral in these
low-frequency moments. It may stabilize or destabilize other modes.

This corollary neither proves superiority over diagonal gains (which already
have a moment-neutral subspace) nor supplies a distributed implementation.
Communication, signal availability, delays, gain bounds and nonlinear
validation would require a new controller contract. No numerical result below
is attributed to this unimplemented architecture.

\subsection{Numerical scope and unresolved publication claim}
The coefficient identities concern the full linearized physical-supply model.
Finite-amplitude nonlinear simulations test operational feasibility separately.
The symmetry-restored germ is explicitly documented; rounding defects in raw
CSV matrices are not silently treated as exact rotational symmetry. The
following values and decisions are generated from the experiment tables.

The tools used here---Laurent residues, response moments, rearrangement and
constrained QP---are established. Delayed PLL tuning also predates this work
\cite{wilson}; signed-area constraints have an established control literature
\cite{aalipourmoments}. A journal-level contribution requires a demonstrated
engineering value for the structural constraint and a fair comparison with
modal tuning that omits it. No maximum replacement, global gain optimum,
robust operating-region certificate, or first-in-literature claim follows
from this section.

% GENERATED_MOMENT_RESULTS_PLACEHOLDER
